\section{Discussion}\label{sec:discussion}

\subsection{Principal findings}
Holding the model, data and evaluation protocol fixed, we found four things. First, the node-feature
representation strongly affects three-class AD/FTD/control classification (partial
$\eta^2=\anovaRepEta$), with wavelet, spectral and fused features leading. Second, how the graph is
built matters little (partial $\eta^2=\anovaEdgeEta$), the best representation does not depend on it,
and the GCN performs no better than a random forest or SVM on the same features. Third, splitting
windows instead of subjects inflates subject-level macro-F1 from about 0.59 to about 0.99 on the
identical pipeline, reproducing the top of the published range. Fourth, the
blood--brain-barrier-associated PSWE marker is more frequent in AD and FTD than in controls, but this
excess is explained by continuous background slowing; defined relative to each person's background,
transient slowing events are, if anything, less frequent in dementia.

\subsection{Representation matters more than topology}
The size and consistency of the representation effect, against a small and non-interacting graph
effect, indicate that the information a GNN can use on this benchmark lives mostly in the node
features. Several observations point the same way. Spatial graphs --- which are identical for every
subject and therefore carry no diagnostic information of their own --- did at least as well as
functional and hybrid graphs. A degree-preserving random graph cost about as much as removing the
graph altogether, which suggests that message passing helps, where it helps at all, by smoothing
features over neighbouring electrodes rather than by exploiting a particular connectivity pattern.
And non-graph classifiers on flattened features matched the GCN cell for cell. These results agree
with the independent conclusion of the 2026 scoping review that apparent representational gains on
this dataset do not survive leakage-free evaluation \cite{R1}, and they temper the case for
graph-based modelling made in recent work \cite{R2,R8,F7}: on 19-electrode scalp EEG from 88
subjects, we see no evidence that graph structure adds information beyond good node features. Whether
this holds with source-space nodes, learned adjacency or much larger cohorts remains open.

\subsection{Neurophysiological interpretation}
Representations that describe each electrode's own spectrum (P1, P3) outperformed one built only from
inter-electrode phase coupling (P2), and the functional graph did not improve on the spatial one.
Within the limits of scalp EEG, this favours local oscillatory slowing over altered large-scale
coupling as the dominant discriminative signal, consistent with spectral slowing being the most
replicated EEG marker of AD \cite{R3,R15} and with the posterior-alpha-to-diffuse-theta shift visible
in our data (Figure~\ref{fig:topo}). It does not contradict network accounts of dementia
\cite{F12,F9}: alpha-band wPLI estimated within \SI{10}{\second} windows is noisy, and electrode-level
graphs are a coarse proxy for cortical networks. The error pattern is clinically recognisable:
controls are identified best, while FTD --- the smallest group, and less slowed than AD --- is
identified worst and is usually mistaken for AD. Separating the two dementias remains the hard part of the problem.

\subsection{What the PSWE results mean for an EEG marker of BBB dysfunction}
With the fixed threshold, our data agree in direction with Milikovsky et al.: PSWE rates were roughly
three times higher in AD than in controls after adjusting for age and sex, and FTD --- not previously
studied --- was also elevated. Two further analyses change the interpretation. The fixed-threshold
rate was almost collinear with relative delta+theta power, and adjusting for that slowing reversed
the AD effect. When events were instead defined as transient drops below each channel's own
background, AD and FTD showed fewer of them than controls, under every variant we examined. Together
these results suggest that, in resting scalp EEG with diffuse slowing, a median power frequency below
\SI{6}{\hertz} mostly flags a slow \emph{background} rather than \emph{paroxysmal} events.

We draw three cautious conclusions. First, a fixed-threshold PSWE rate should not be read as evidence
of BBB dysfunction without controlling for background slowing. Second, background-relative
definitions are needed if PSWEs are to capture transient cortical dysfunction, although they bring
their own floor effect in very slow recordings, which we observed in the Hz-drop variants. Third, our
data cannot test the BBB link itself: ds004504 contains no BBB imaging, the original detection
details include parameters we could not verify from the full text, and Milikovsky et al.\ studied
different patients under different recording conditions. The decisive test is a cohort with both EEG
and dynamic contrast-enhanced MRI.

\subsection{Implications for how results on this benchmark are read}
On one pipeline, the choice between subject-level and window-level splitting moved subject-level
macro-F1 by about 0.4 --- roughly three times the gap between the best and worst representation, and
several times larger than any graph or architecture change we tested. Figures in the high 90s on ds004504 should therefore be read alongside their
splitting protocol, and studies that report window-level metrics or allow a subject's windows into
both training and test sets cannot be compared with subject-level results. We name no individual
study: the point is that the protocol, not the model, determines most of the reported number, as
argued independently in \cite{R1,F8}. The permutation test confirms that the honest result is
nonetheless real.

\subsection{Limitations}
The cohort is small: test folds contain 17--18 subjects, and under the conservative Nadeau--Bengio
correction no pairwise difference between representations, and no single ablation, reached
significance. We report both corrected and uncorrected tests and give effect sizes, but small
differences between leading representations should not be over-interpreted. All data come from one
site, device and protocol, and FTD is the smallest class ($n=23$). The planned external validation on
an independent cohort (ds004584) and two planned ablations (window length and ICA on/off) were not
run for this draft. Hyperparameters were fixed a priori and not tuned; early stopping on a
validation set of about 14 subjects stopped training early, which may disadvantage higher-capacity
configurations such as raw windows (P5, about 84k parameters). Calibration was poor. Graphs are
defined on electrodes rather than cortical sources. The demographic floor (\demoFOne) shows residual
age and sex information, which residualisation reduced but did not remove. For PSWEs, the MPF
frequency range was not reported in the source and was set by us, we detected events per channel
rather than over averaged regions, the background-relative analysis is post hoc, and no BBB measure
is available.

\subsection{Future work}
The most valuable next study is a multi-site replication of this exact protocol, since a single-site
design cannot establish device or population invariance. Natural extensions are source-space graphs
\cite{F7}, learned adjacency \cite{F4}, spatio-temporal graphs over window sequences \cite{R10,F6},
an architecture sweep under a fixed representation (the transpose of this design) \cite{F5}, an
MCI class, and --- for the BBB question --- simultaneous EEG and dynamic contrast-enhanced MRI to test
whether any PSWE definition tracks measured barrier leakage \cite{B5,B8}.
